\documentclass{article}
%\usepackage[spanish]{babel}
\usepackage[pdftex]{graphicx,color} 
\usepackage[utf8]{inputenc}
\usepackage{listings}
\usepackage{colortbl}
\usepackage{textcomp}
\usepackage{listings}
\usepackage{color}
\usepackage{wrapfig}
\usepackage{amsfonts}
\usepackage{amssymb}
\usepackage{algorithmic}
\usepackage{algorithm}
\usepackage[export]{adjustbox}
%\usepackage{amslatex}
\definecolor{lbcolor}{rgb}{0.9,0.9,0.9}
\lstset{
        backgroundcolor=, %\color{lbcolor},
        tabsize=4,
        rulecolor=,
        language=C,
        basicstyle=\scriptsize,
        upquote=true,
        aboveskip={1.5\baselineskip},
        %columns=fixed,
        showstringspaces=false,
        extendedchars=true,
        breaklines=true,
        prebreak = \raisebox{0ex}[0ex][0ex]{\ensuremath{\hookleftarrow}},
        frame=single,
        showtabs=false,
        showspaces=false,
        showstringspaces=false,
        identifierstyle=\ttfamily,
        keywordstyle=\color[rgb]{0,0,1},
        commentstyle=\color[rgb]{0.133,0.545,0.133},
        stringstyle=\color[rgb]{0.627,0.126,0.941},
}

\topmargin -0.5in
\textwidth 6.8in
\textheight 9.0in
\oddsidemargin 0.0in
\evensidemargin 0.0in

\usepackage{enumitem}
\usepackage{multirow}
\usepackage{hhline}
\usepackage{multicol}
%ointpoints{punto}{puntos}

\begin{document}
\vspace{.5pc}
\centerline{\sc Clase Activa 1}
\centerline{\sc ELO320 - Estructura de Datos y Algoritmos}
\centerline{\sc Prof. Nicolás Gálvez Ramírez}
\vspace{2pc}
%\noindent {\bf Instrucciones}:
%\begin{itemize}
%    \item Conteste cada pregunta en una hoja separada. Éstas serán entregadas de tal forma al finalizar el certamen.
%    \item Si decide no contestar una pregunta, debe entregar una hoja sin respuestas en su lugar. 
%    \item Rotule claramente cada una de sus hojas con su Nombre Completo, Rol USM y Número de Paralelo.
%    \item El certamen tiene una duración de 60 minutos.
%\end{itemize}

%\noindent {\bf Preguntas}
%\medskip
\begin{itemize}

\begin{minipage}[h]{0.45\textwidth}
\item \textbf{Operadores Binarios}:

\begin{tabular}{|l|l|l|l|}\hline 
	\textbf{Operador} &\textbf{C} & \textbf{Python}   \\ \hline \hline
	Negación lógica & !  &  $\mbox{not}$ \\  \hline
	Dirección       & \&  &  \emph{No existe} \\  \hline
	Derreferencia   & * &  \emph{No existe} \\  \hline
%	Complemento a 1 & $\sim$ & $\sim$  \\  \hline
	Incremento      & ++ &  \emph{No existe} \\  \hline
	Decremento      & - - & \emph{No existe} \\  \hline
	Signo más       & + &  + \\  \hline
	Negativo        & - &  - \\  \hline
	Casting         & (tipo) &  \emph{No existe} \\  \hline
\end{tabular}\\[1em]
Los operadores unarios agrupan de derecha a izquierda. Ejemplo: \textcolor{blue}{$-+ x$} equivalente a \textcolor{blue}{$-(+ x)$}.
\end{minipage}
\hspace{0.5cm}
\begin{minipage}[h]{0.45\textwidth}
\item \textbf{Operadores Binarios}:

\begin{tabular}{|l|l|l|l|}\hline
	\textbf{Operador} &\textbf{C} & \textbf{Python}   \\ \hline \hline
	Suma         & +  & +  \\  \hline
	Resta       & - &  - \\  \hline
	Multiplicación & * & *  \\  \hline
	División      & / &  / \\  \hline
	División entera & /  & //  \\  \hline
	Módulo      & \% & \% \\  \hline
	Exponenciación   & \emph{No existe}  &  ** \\  \hline
\end{tabular}\\[1em]

\textbf{División entera en C según tipo de dato.} Ejemplo, $\frac{5}{4} = 1$ y $\frac{5.0}{4} = 1.25 $. 
\end{minipage}

\item \textbf{Precedencia de operadores}:

\begin{center}
\begin{tabular}{|l|l||l|l|c|}\hline 
	\multirow{2}{*}{\textbf{Precedencia}}  & \multirow{2}{*}{\textbf{Descripción}} & 
	\multicolumn{2}{|c|}{\textbf{Operador}} & \multirow{2}{*}{\textbf{Asociatividad}} \\ \hhline{~~--~}
	& & \textbf{C} & \textbf{Python} &  \\ \hline \hline
	alta & exponente &  & $**$ & $\leftarrow$  \\	\hline
	& unarios (postfijos) & \tiny $++$ $--$ &  & $\rightarrow$ \\	\hline
	& unarios (prefijos) & \tiny$++$ $--$ $\sim$ ! * \& + - & \tiny $\sim$ + - & $\leftarrow$ \\	\hline
	& casting de tipo &(type) & & $\leftarrow$ \\	\hline
	& multiplicativos & \tiny * / \% & \tiny * / // \%	 & $\rightarrow$  \\	\hline
	& aditivos & \tiny $+$ $-$ & \tiny $+$ $-$ & $\rightarrow$  \\	\hline
%	& corrimiento & \tiny $<<$ $>>$ & \tiny $<<$ $>>$ & $\rightarrow$  \\	\hline
%	& bitwise AND  & \& & \&    & $\rightarrow$  \\	\hline
%	& bitwise XOR & $\wedge$ & $\wedge$ & $\rightarrow$  \\	\hline
%	& bitwise OR  & $|$ & $|$ & $\rightarrow$  \\	\hline
	baja & NOT lógico  &	& $\mathbf{not}$ & $\leftarrow$ \\	\hline
\end{tabular}
\end{center}

\bigskip
\begin{minipage}[h]{0.45\textwidth}
\item \textbf{Entender coherciones}:
\begin{lstlisting}[language=C, basicstyle=\scriptsize, escapeinside={|}{|}]
float   fVal;
double  dVal;
int   iVal;
unsigned long ulVal;

dVal = iVal * ulVal; //iVal es convertido a unsigned long.
// El resultado es convertido a double 

dVal = ulVal + fVal; //ulVal convertido a float.
//El resultado es convertido a double            
	\end{lstlisting}
\end{minipage}
\hspace{0.5cm}
\begin{minipage}[h]{0.45\textwidth}
\item \textbf{Asignaciones}:

\medskip
\begin{tabular}{|l|l|l|}\hline
	\textbf{Operador} &\textbf{C} & \textbf{Python}   \\ \hline 
	Asignación         & $=$  & $=$  \\  \hline
	Suma y asignación       & $+=$ & $+=$ \\  \hline
	Resta y asignación          & $-=$ & $-=$  \\  \hline
	Multiplicación y asignación       & $*=$ & $*=$ \\  \hline
	División y asignación   & $/=$ & $/=$  \\  \hline
	Módulo y asignación   & $\%=$ & $\%=$  \\  \hline
%	Bitwise AND y asignación   &  $\&=$ &  \\  \hline
%	Bitwise OR y asignación   & $|=$ &   \\  \hline
%	Bitwise XOR y asignación   & $\wedge=$ &  \\  \hline
%	Corrimiento der. y asignación   & $>>=$ &   \\  \hline
%	Corrimiento izq. y asignación   & $<<=$ &   \\  \hline
	Incremento    & $++$ &  \\  \hline
	Decremento   & $--$ &   \\  \hline
	División entera y asignación   &  & $//=$  \\  \hline
	Exponente y asignación   & & $**=$  \\  \hline
\end{tabular}
\end{minipage}

\item \textbf{Unarios de Asignación - C:}

\begin{minipage}[h]{0.45\textwidth}
\begin{lstlisting}[language=C, basicstyle=\scriptsize] 
//pre-incremental
sum = ++contador; 

//equivale a:

contador = contador + 1;
sum = contador;
\end{lstlisting}
\end{minipage}
\begin{minipage}[h]{0.45\textwidth}
\begin{lstlisting}[language=C, basicstyle=\scriptsize] 
//post-incremental
sum = contador++;
 
//equivale a:
aux = contador;
contador = contador + 1;
sum = aux;
\end{lstlisting}
\end{minipage}
\newpage
\item Operador Ternario:

\begin{minipage}[h]{0.45\textwidth}
\begin{lstlisting}[language=C, basicstyle=\scriptsize] 
retorno = op1 ? op2 : op3;  
\end{lstlisting}
{\scriptsize \textbf{Si op1 es verdadero, asigna op2; sino, asigna op3.}}\\
\end{minipage}
\hspace{0.5cm}
\begin{minipage}[h]{0.45\textwidth}
\begin{lstlisting}[language=C, basicstyle=\scriptsize] 
cuenta = (a > 0) ? 0 : 1;

//Equivale a:
if (a > 0)
	cuenta = 0;
else cuenta = 1;
\end{lstlisting}
\end{minipage}

\begin{minipage}[h]{0.45\textwidth}
\item \textbf{Expresiones Relacionales y Lógicas}:

\begin{tabular}{|l|l|l|}\hline \hline
	\multicolumn{3}{|c|}{\textbf{Operadores Relacionales}} \\ \hline
	\textbf{Operador} &\textbf{C} & \textbf{Python}   \\ \hline 
	Igual         & $==$  & $==$  \\  \hline
	Distinto       & $!=$ & $!=$ $<>$ \\  \hline
	Mayor          & $>$ & $>$  \\  \hline
	Menor       & $<$ & $<$ \\  \hline
	Mayor o igual   & $>=$ & $>=$  \\  \hline
	Menor o igual   & $<=$ & $<=$  \\  \hline
	\multicolumn{3}{|c|}{\textbf{Operadores Booleanos}} \\ \hline
	\textbf{Operador} &\textbf{C}  & \textbf{Python}   \\ \hline 
	NOT   & !  &   \textbf{not} \\  \hline
	AND   & $\&\&$  & \textbf{and}  \\  \hline
	OR   & $||$  & \textbf{or}  \\  \hline
\end{tabular}
\end{minipage}
\hspace{0.5cm}
\begin{minipage}[h]{0.45\textwidth}
\item \textbf{if-else if-else}:

\begin{lstlisting}[language=C, basicstyle=\scriptsize, escapeinside={|}{|}] 
if(condicion_1) {
    //....
} else if(condicion_k) {
    //....
} else {
    //....
}
\end{lstlisting}
\end{minipage}

\begin{minipage}[h]{0.45\textwidth}

\item \textbf{switch}:
\begin{lstlisting}[language=C, basicstyle=\scriptsize, escapeinside={|}{|}] 
switch(x){ //Debe ser una variable
case x1:                                     
	//.... 
	break; 
case x2: 
	//.... 
	break; 
//.... 
case xN: 
	//.... 
	break;  
default: //opcional
	//....
} 
\end{lstlisting}
\end{minipage}
\hspace{0.5cm}
\begin{minipage}[h]{0.45\textwidth}
\item \textbf{while - do while}:
\begin{lstlisting}[language=C, basicstyle=\scriptsize, escapeinside={|}{|}] 
while(condicion) { 
    //...
}

do {
    //...
} while(condicion);
		\end{lstlisting}

\end{minipage}

\begin{minipage}[h]{0.45\textwidth}
\item \textbf{for}:
\begin{lstlisting}[language=C, basicstyle=\scriptsize, escapeinside={|}{|}] 
//En general:
for(sentencia_inicial;condicion;sentencia_fin_ciclo) { 
    //...
} 

//Uso corriente
for(i=0; i<N; i++) { 
    //...
} 

//Equivalente a:
i=0;
while(i<N) {
    //...
    i++;
}
\end{lstlisting}
\end{minipage}
\hspace{0.5cm}
\begin{minipage}[h]{0.45\textwidth}
\item Corto Circuito o Término anticipado de la evaluación: 

\begin{lstlisting}[language=C, basicstyle=\scriptsize] 
//Sin tildes por codificacion de caracteres.
//Si primera condicion se cumple, la expresion es true, no se evalua el resto
if((a >= 0) || (b < 10))

//Si primera condicion no se cumple, la expresion es false, no se evalua el resto.
while ( c > 0 && c+=1 < 10 ))
\end{lstlisting}
\end{minipage}
\end{itemize}
\end{document}

